% GENERATED FILE. Edit canonical lessons, then run npm run generate:books.
% canonical-source-hash: fnv1a64:9ba22da30e28d61c
% canonical-lessons: HI-C399-lain, HI-C399-injan, HI-C399-slipar, HI-C399-siti, HI-C399-sarakna

\chapter{Lines and Carriages}
\label{ch:hi-lines-and-carriages}
\hlchaptermodality{399}

\begin{chapteropening}
\textbf{By the end of this chapter:} \emph{I can talk about the railway line, the engine and the carriages, and a train sliding forward.}

Complete the last lesson of chapter 399: I can talk about the railway line, the engine and the carriages, and a train sliding forward.
\end{chapteropening}

\section[\texorpdfstring{\hi{लाइन} (lāin)}{लाइन (lāin)}]{\hi{लाइन} (lāin) — a line; a railway line}

\label{lesson:HI-C399-lain}



\begin{quote}
Before the new one: say the Hindi for to load (onto), then the Hindi for to seat (someone).
\end{quote}

\subsection*{You'll want to know: \hi{लाइन}}
\textbf{\hi{लाइन}} — \emph{lāin} — \textquotedblleft{}a line; a railway line\textquotedblright{}.

\begin{quote}
\hi{रेल} \hi{की} \hi{लाइन} \hi{नदी} \hi{के} \hi{किनारे} \hi{जाती} \hi{है।}
\end{quote}

\begin{quote}
\emph{rel kī lāin nadī ke kināre jātī hai.}
\end{quote}

\begin{quote}
The railway line runs along the river bank.
\end{quote}

Feminine.

Also a queue of people: \hi{लाइन} \hi{में} \hi{खड़े} \hi{रहिए}, stay in the line.

\begin{cousinweb}
\textbf{\hi{लाइन}} (\emph{lāin}) is borrowed from English \emph{line}, used in Hindi for a railway line, a row of people and a line on a page.
\end{cousinweb}

\begin{grammarlens}[title={What you've built}]
One more word for getting around.
\end{grammarlens}

\subsection*{Your turn}
\begin{itemize}
\raggedright
  \item \emph{Say it:} \emph{lāin}
  \item \emph{Say it:} \emph{lāin}, once more
  \item \emph{Say it:} \emph{rel kī lāin nadī ke kināre jātī hai.} — The railway line runs along the river bank.
  \item \emph{Recall:} say the Hindi for to load (onto), then the Hindi for to seat (someone), then say \emph{lāin} again
\end{itemize}

\begin{tcolorbox}[breakable,colback=teal!4,colframe=teal!35!black,title={Before you move on}]
What does \hi{लाइन} mean? (\textquotedblleft{}A line; a railway line\textquotedblright{}.) Say it once more.
\end{tcolorbox}
\section[\texorpdfstring{\hi{इंजन} (injan)}{इंजन (injan)}]{\hi{इंजन} (injan) — an engine}

\label{lesson:HI-C399-injan}



\begin{quote}
Before the new one: say the Hindi for to seat (someone), then the Hindi for a line.
\end{quote}

\subsection*{You'll want to know: \hi{इंजन}}
\textbf{\hi{इंजन}} — \emph{injan} — \textquotedblleft{}an engine\textquotedblright{}.

\begin{quote}
\hi{इंजन} \hi{बहुत} \hi{पुराना} \hi{था}, \hi{पर} \hi{तेज़} \hi{चलता} \hi{था।}
\end{quote}

\begin{quote}
\emph{injan bahut purānā thā, par tez caltā thā.}
\end{quote}

\begin{quote}
The engine was very old, but it ran fast.
\end{quote}

Masculine.

\begin{cousinweb}
\textbf{\hi{इंजन}} (\emph{injan}) is English \emph{engine}, taken over with the railways; the English soft g is said with \hi{ज}.
\end{cousinweb}

\begin{grammarlens}[title={What you've built}]
One more word for getting around.
\end{grammarlens}

\subsection*{Your turn}
\begin{itemize}
\raggedright
  \item \emph{Say it:} \emph{injan}
  \item \emph{Say it:} \emph{injan}, once more
  \item \emph{Say it:} \emph{injan bahut purānā thā, par tez caltā thā.} — The engine was very old, but it ran fast.
  \item \emph{Recall:} say the Hindi for to seat (someone), then the Hindi for a line, then say \emph{injan} again
\end{itemize}

\begin{tcolorbox}[breakable,colback=teal!4,colframe=teal!35!black,title={Before you move on}]
What does \hi{इंजन} mean? (\textquotedblleft{}An engine\textquotedblright{}.) Say it once more.
\end{tcolorbox}
\section[\texorpdfstring{\hi{स्लीपर} (slīpar)}{स्लीपर (slīpar)}]{\hi{स्लीपर} (slīpar) — a sleeper carriage; sleeper class}

\label{lesson:HI-C399-slipar}



\begin{quote}
Before the new one: say the Hindi for a line, then the Hindi for an engine.
\end{quote}

\subsection*{You'll want to know: \hi{स्लीपर}}
\textbf{\hi{स्लीपर}} — \emph{slīpar} — \textquotedblleft{}a sleeper carriage; sleeper class\textquotedblright{}.

\begin{quote}
\hi{मेरा} \hi{टिकट} \hi{स्लीपर} \hi{में} \hi{है।}
\end{quote}

\begin{quote}
\emph{merā ṭikaṭ slīpar meṁ hai.}
\end{quote}

\begin{quote}
My ticket is for the sleeper class.
\end{quote}

Masculine.

\begin{cousinweb}
\textbf{\hi{स्लीपर}} (\emph{slīpar}) is borrowed from English \emph{sleeper}, short for a sleeper carriage, where the seats fold out into berths for the night.
\end{cousinweb}

\begin{grammarlens}[title={What you've built}]
One more word for getting around.
\end{grammarlens}

\subsection*{Your turn}
\begin{itemize}
\raggedright
  \item \emph{Say it:} \emph{slīpar}
  \item \emph{Say it:} \emph{slīpar}, once more
  \item \emph{Say it:} \emph{merā ṭikaṭ slīpar meṁ hai.} — My ticket is for the sleeper class.
  \item \emph{Recall:} say the Hindi for a line, then the Hindi for an engine, then say \emph{slīpar} again
\end{itemize}

\begin{tcolorbox}[breakable,colback=teal!4,colframe=teal!35!black,title={Before you move on}]
What does \hi{स्लीपर} mean? (\textquotedblleft{}A sleeper carriage; sleeper class\textquotedblright{}.) Say it once more.
\end{tcolorbox}
\section[\texorpdfstring{\hi{सीटी} (sīṭī)}{सीटी (sīṭī)}]{\hi{सीटी} (sīṭī) — a whistle}

\label{lesson:HI-C399-siti}



\begin{quote}
Before the new one: say the Hindi for an engine, then the Hindi for a sleeper carriage.
\end{quote}

\subsection*{You'll want to know: \hi{सीटी}}
\textbf{\hi{सीटी}} — \emph{sīṭī} — \textquotedblleft{}a whistle\textquotedblright{}.

\begin{quote}
\hi{ट्रेन} \hi{ने} \hi{लंबी} \hi{सीटी} \hi{दी।}
\end{quote}

\begin{quote}
\emph{ṭren ne lambī sīṭī dī.}
\end{quote}

\begin{quote}
The train gave a long whistle.
\end{quote}

Feminine.

\hi{सीटी} \hi{देना} or \hi{सीटी} \hi{बजाना}, to blow a whistle.

\begin{cousinweb}
\textbf{\hi{सीटी}} (\emph{sīṭī}) goes back to \emph{*sīṭṭa}, \textquotedblleft{}whistle\textquotedblright{}, a form CDIAL reconstructs from its descendants and treats as imitative of the sound itself (CDIAL 13427).
\end{cousinweb}

\begin{grammarlens}[title={What you've built}]
One more word for getting around.
\end{grammarlens}

\subsection*{Your turn}
\begin{itemize}
\raggedright
  \item \emph{Say it:} \emph{sīṭī}
  \item \emph{Say it:} \emph{sīṭī}, once more
  \item \emph{Say it:} \emph{ṭren ne lambī sīṭī dī.} — The train gave a long whistle.
  \item \emph{Recall:} say the Hindi for an engine, then the Hindi for a sleeper carriage, then say \emph{sīṭī} again
\end{itemize}

\begin{tcolorbox}[breakable,colback=teal!4,colframe=teal!35!black,title={Before you move on}]
What does \hi{सीटी} mean? (\textquotedblleft{}A whistle\textquotedblright{}.) Say it once more.
\end{tcolorbox}
\section[\texorpdfstring{\hi{सरकना} (saraknā)}{सरकना (saraknā)}]{\hi{सरकना} (saraknā) — to slide; to shift along}

\label{lesson:HI-C399-sarakna}



\begin{quote}
Before the new one: say the Hindi for a sleeper carriage, then the Hindi for a whistle.
\end{quote}

\subsection*{You'll want to know: \hi{सरकना}}
\textbf{\hi{सरकना}} — \emph{saraknā} — \textquotedblleft{}to slide; to shift along\textquotedblright{}.

\begin{quote}
\hi{ट्रेन} \hi{धीरे} \hi{आगे} \hi{सरक} \hi{रही} \hi{थी।}
\end{quote}

\begin{quote}
\emph{ṭren dhīre āge sarak rahī thī.}
\end{quote}

\begin{quote}
The train was sliding slowly forward.
\end{quote}

Past (it …): \hi{वह} \hi{मेज़} \hi{थोड़ी} \hi{आगे} \hi{सरकी}, \emph{vah mez thoṛī āge sarkī}.

\begin{cousinweb}
\textbf{\hi{सरकना}} (\emph{saraknā}) goes back to Sanskrit \emph{sarati}, \textquotedblleft{}runs, flows, glides\textquotedblright{} (CDIAL 13250), which lists Hindi \emph{saraknā} in the sense of drawing back; in everyday Hindi it is a slow slide along a surface.
\end{cousinweb}

\begin{grammarlens}[title={What you've built}]
One more word for getting around.
\end{grammarlens}

\subsection*{Your turn}
\begin{itemize}
\raggedright
  \item \emph{Say it:} \emph{saraknā}
  \item \emph{Say it:} \emph{saraknā}, once more
  \item \emph{Say it:} \emph{ṭren dhīre āge sarak rahī thī.} — The train was sliding slowly forward.
  \item \emph{Recall:} say the Hindi for a sleeper carriage, then the Hindi for a whistle, then say \emph{saraknā} again
\end{itemize}

\begin{tcolorbox}[breakable,colback=teal!4,colframe=teal!35!black,title={Before you move on}]
What does \hi{सरकना} mean? (\textquotedblleft{}To slide; to shift along\textquotedblright{}.) Say it once more.
\end{tcolorbox}
